\SourceLeaf{89}%

If we now suppose that the functions $fx$ and $\varphi x$ are related in such a way that

\[ \Phi x = \psi x \mathord{.} Fx \]

\noindent or, equivalently, that

\[ \varphi x = \int \psi x \mathord{.} \frac{dfx}{dx}\,dx, \]

\noindent and consider that

\[ \frac{dfx}{du} = \psi x\,\frac{dfx}{dt}, \]

\noindent then

\[ \frac{dfx}{du} = \frac{d\varphi x}{dt}. \]

\noindent Hence [it follows], since

\[ \frac{d\varphi t}{dt} = \Phi t = \psi t \mathord{.} Ft, \]

\[\begin{aligned}
T_1 + 2 T_2 u + 3 T_3 uu + \text{u.\,s.\,w.}
&= \psi t\mathord{.}Ft+\frac{d\Theta_1}{dt}u\\
&\quad+\frac{d\Theta_2}{dt}uu+\frac{d\Theta_3}{dt}u^3+\text{u.\,s.\,w.}
\end{aligned}\]

\Needspace{4\baselineskip}
\noindent We therefore have first of all

\[ T_1 = \psi t \mathord{.} Ft. \]

\noindent Thus, by the preceding remark, $T_1$ must become $\Theta_1$ if in the expression for the former function one replaces $Ft$ by $\Phi t$, i.\,e. $\psi t \mathord{.} Ft$. This gives

\[ \Theta_1 = (\psi t)^2 \mathord{.} Ft. \]

\noindent From this we further obtain

\[ T_2 = \frac{1}{2}\,\frac{d(\psi t)^2 Ft}{dt}, \]

\noindent and then, again by the preceding remark,

\[ \Theta_2 = \frac{1}{2}\,\frac{d(\psi t)^3 Ft}{dt}. \]

\noindent Then

\[ T_3 = \frac{1}{2.3}\,\frac{dd(\psi t)^3 Ft}{dt^2}, \]
\[ \Theta_3 = \frac{1}{2.3}\,\frac{dd(\psi t)^4 Ft}{dt^2}, \]
\[ T_4 = \frac{1}{2.3.4}\,\frac{d^3(\psi t)^4 Ft}{dt^3}, \]

\begin{center}and so on.\end{center}
